% ============================================================
% Capitolo 8 — Stati temporizzati (20260511 Stati temporizzati2.pdf)
% ============================================================
\chapter{Stati temporizzati}
\label{cap:stati-temporizzati}

\section{Introduzione}
Il timer TON (Capitolo~\ref{cap:timer-ton}) viene utilizzato attraverso il nuovo concetto di \textbf{stato temporizzato}:
\begin{itemize}
    \item uno stato temporizzato è tale che il controllo rimane in esso per \textbf{al più un tempo fissato};
    \item scaduto questo tempo, una \textbf{transizione autonoma} (detta di \textbf{timeout}) porta il controllo verso un altro stato;
    \item altre transizioni possono portare il controllo fuori dallo stato \textbf{prima} del timeout.
\end{itemize}

\section{Simbolismo grafico}
\begin{enumerate}
    \item Lo stato temporizzato presenta un \textbf{riquadro} che indica il tempo massimo $T$ in cui il controllo rimarrà in esso;
    \item la transizione di \textbf{timeout} presenta una $T$ e un pallino, indicandone la natura temporale e autonoma.
\end{enumerate}

\begin{figure}[htbp]
\centering
\begin{tikzpicture}[>=stealth, every node/.style={font=\small}]
    \node[state, draw, minimum size=1.6cm, label={[anchor=south east, font=\footnotesize, draw, inner sep=1pt]north east:$T$}] (s1) {$S_1$};
    \node[state, draw, minimum size=1.6cm, right=4.5cm of s1] (s2) {$S_2$};

    \draw[->] (s1) edge[bend left=20] node[above] {$t_1$: timeout ($T$, $\bullet$)} (s2);
    \draw[->] (s1) edge[bend right=35] node[below] {$t_2$: evento \texttt{ev} (trigger)} (s2);
\end{tikzpicture}
\caption{Notazione dello stato temporizzato: se l'evento \texttt{ev} che agisce da trigger giunge prima dello scadere del tempo $T$, verrà intrapresa la transizione $t_2$; altrimenti, allo scadere di $T$, scatterà la transizione autonoma di timeout $t_1$.}
\label{fig:stato-temporizzato}
\end{figure}

\section{Esempio: semaforo}
Un primo esempio di stato temporizzato è la realizzazione del semaforo stradale: la transizione da giallo (\texttt{Y}) a rosso (\texttt{R}) è sempre autonoma, ma ora \emph{sappiamo a che tempo avviene} ($T = 5\,s$):
\begin{lstlisting}[basicstyle=\ttfamily\small]
Stati_Sem.Y:
  timer(IN := TRUE, PT := T#5S );
  IF timer.Q THEN
    stato := Stati_Sem.R;
    timer(IN := FALSE);
  END_IF
\end{lstlisting}

\section{Spia lampeggiante}
Una spia lampeggiante passa dallo stato acceso (\texttt{On}) allo stato spento (\texttt{Off}), ciascuno dei due stati ha una durata predefinita: il comportamento globale è dato da \textbf{due stati temporizzati collegati tra loro}.

\subsection{Implementazione (versione base)}
\begin{lstlisting}[basicstyle=\ttfamily\small]
FUNCTION_BLOCK Light
VAR_OUTPUT
  stato : Light_St := Light_St.Init;
END_VAR
VAR
  timer : TON;
END_VAR

CASE stato OF
  Light_St.Init:
    stato := Light_St.On;

  Light_St.Off:
    timer(IN:=TRUE, PT := T#1S);
    IF timer.Q THEN
      stato := Light_St.On;
      timer(IN := FALSE );
    END_IF

  Light_St.On:
    timer(IN:=TRUE, PT := T#2S);
    IF timer.Q THEN
      stato := Light_St.Off;
      timer(IN := FALSE );
    END_IF
END_CASE
\end{lstlisting}

\begin{note}[Regole implementative dello stato temporizzato]
\begin{itemize}
    \item Uno stato temporizzato \textbf{richiama il timer all'inizio del proprio codice} (ad ogni ciclo di scan);
    \item la transizione di timeout è attivata dal segnale logico \texttt{Q} del timer: \texttt{IF timer.Q THEN ...};
    \item \textbf{ogni transizione in uscita} (compresa quella di timeout) \textbf{termina con il reset del timer}: \texttt{timer(IN := FALSE)}.
\end{itemize}
\end{note}

\subsection{Versione 2: stop e restart}
In questa seconda versione la spia lampeggiante può essere \textbf{fermata in ogni momento} con l'evento \texttt{stop} e fatta \textbf{ripartire} con l'evento \texttt{restart}. Ogni stato temporizzato contiene, oltre alla logica di timeout, un \texttt{IF} che verifica l'evento di stop (che consuma l'evento e resetta il timer):
\begin{lstlisting}[basicstyle=\ttfamily\small]
FUNCTION_BLOCK Light
VAR_INPUT
  stop_ev, restart_ev : BOOL := FALSE;
END_VAR
VAR_OUTPUT
  stato : Light_St := Light_St.Init;
END_VAR
VAR
  timer : TON;
END_VAR

CASE stato OF
  Light_St.Init:
    stato := Light_St.On;

  Light_St.Off:
    timer(IN:=TRUE, PT := T#2S);
    IF timer.Q THEN
      stato := Light_St.On;
      timer(IN := FALSE );
    END_IF
    IF stop_ev THEN
      stato := Light_St.Stop;
      stop_ev := FALSE;
      timer(IN := FALSE);
    END_IF

  Light_St.On:
    timer(IN:=TRUE, PT := T#1S);
    IF timer.Q THEN
      stato := Light_St.Off;
      timer(IN := FALSE );
    END_IF
    IF stop_ev THEN
      stato := Light_St.Stop;
      stop_ev := FALSE;
      timer(IN := FALSE);
    END_IF

  Light_St.Stop:
    IF restart_ev THEN
      stato := Light_St.On;
      restart_ev := FALSE;
    END_IF
END_CASE
\end{lstlisting}

\section{Controller semaforico temporizzato}
\subsection{Direttrici abilitate a tempo}
Nel controller semaforico visto in precedenza le due direttrici venivano attivate da due eventi trigger. In questa variante le due direttrici \texttt{ns} (nord--sud) ed \texttt{ew} (est--ovest) sono abilitate \textbf{esattamente per un tempo stabilito} (stati temporizzati).

\subsection{Farm Road}
\begin{defn}[Farm Road]
Strada secondaria che si immette su una strada principale (\emph{main road}). La farm road rimane abilitata \textbf{a richiesta} (evento \texttt{car}) esattamente per un tempo $T = 30\,s$.
\end{defn}

\subsection{Farm Road — Extended Time}
Variante: la farm road rimane abilitata a richiesta (evento \texttt{car}) esattamente per un tempo $T$; \textbf{mentre} la farm road è abilitata, un secondo evento \texttt{car} può \textbf{estendere} il tempo di abilitazione di ulteriori 30 secondi \textbf{ma solo una volta}.

\begin{intuition}[Perché gli stati temporizzati]
Gli stati temporizzati rendono \emph{esplicito nel modello} il tempo di permanenza in uno stato, integrando il timer TON nel formalismo a stati: il timeout è una transizione autonoma (scatta da sola, senza trigger esterno), mentre le altre transizioni possono anticiparlo. In ST si implementano con il pattern: chiamata del timer in testa allo stato, test di \texttt{timer.Q} per il timeout, reset del timer su ogni uscita, e gestione degli eventi con consumo del segnale.
\end{intuition}
